\section*{Sets and Indices}

There are no nontrivial index sets in the implemented model. The model has two product activities:
\[
i \in \{A,B\}.
\]

\section*{Parameters}

From the data file:
\[
p_A = 3
\qquad\text{(profit of product A, code: \texttt{profit\_A})},
\]
\[
p_B = 5
\qquad\text{(profit of product B, code: \texttt{profit\_B})},
\]
\[
t_A = 12
\qquad\text{(assembly time per unit of A in minutes, code: \texttt{assembly\_time\_A})},
\]
\[
t_B = 25
\qquad\text{(assembly time per unit of B in minutes, code: \texttt{assembly\_time\_B})}.
\]

Weekly machine time is read in hours and converted in the code to minutes:
\[
H = 30
\qquad\text{(machine working time in hours, code: \texttt{machine\_working\_time\_hours})},
\]
\[
T = 60H = 1800
\qquad\text{(machine working time in minutes, code: \texttt{machine\_working\_time\_minutes})}.
\]

The production-ratio data value is ``\(5:2\)'', implemented as constants
\[
r_A = 5
\qquad\text{(code: \texttt{ratio\_A})},
\]
\[
r_B = 2
\qquad\text{(code: \texttt{ratio\_B})},
\]
with the intended requirement that at least \(2\) units of \(B\) are produced for every \(5\) units of \(A\).

\section*{Decision Variables}

\[
x_A \ge 0
\qquad\text{(production quantity of product A, code: \texttt{A})},
\]
\[
x_B \ge 0
\qquad\text{(production quantity of product B, code: \texttt{B})}.
\]

Both variables are continuous in the implementation.

\section*{Objective}

Maximize total weekly profit:
\[
\max \; z = p_A x_A + p_B x_B.
\]

With the given data, this is
\[
\max \; z = 3x_A + 5x_B.
\]

\section*{Constraints}

Machine-time capacity:
\[
t_A x_A + t_B x_B \le T.
\]

With the given data,
\[
12x_A + 25x_B \le 1800.
\]

Production-ratio constraint as implemented:
\[
r_B x_A - r_A x_B \le 0.
\]

Equivalently,
\[
x_B \ge \frac{r_B}{r_A}x_A.
\]

With the given data,
\[
2x_A - 5x_B \le 0,
\]
equivalently
\[
x_B \ge \frac{2}{5}x_A.
\]

Nonnegativity:
\[
x_A \ge 0,\qquad x_B \ge 0.
\]

\section*{Notes on Implementation}

The source code formulates and solves a linear program using PuLP with the Gurobi command-line solver interface (\texttt{pulp.GUROBI\_CMD(msg=0)}).

The variables are explicitly declared as continuous and nonnegative; there are no integrality restrictions.

The ratio parameter from the CSV is parsed by the utility code as a tuple, but the optimization model itself hard-codes
\[
\texttt{ratio\_A}=5,\qquad \texttt{ratio\_B}=2,
\]
and imposes the constraint
\[
\texttt{ratio\_B} \cdot A - \texttt{ratio\_A} \cdot B \le 0.
\]

The machine time in the CSV is provided in hours and converted in the code to minutes before building the capacity constraint.
